\section{Related Work}
\label{sec:related}
\subsection{Efficient GUI Agents}
GUI agents ground actions in an accumulating stream of high-resolution screenshots~\citep{hong2024cogagent,guiowl2025}, inflating both computation and memory.
%
Existing methods fall into \textit{training-based designs} and \textit{training-free designs}.
%
\textbf{Training-based Designs.}
This paradigm learns cheaper perception by retraining the architecture or the history policy.
%
For example, CogAgent~\citep{hong2024cogagent} couples a low-resolution backbone with a high-resolution cross-attention module.
%
ReVision~\citep{abaskohi2026revision} learns to cut temporal visual redundancy along the trajectory.
%
While effective, these designs demand extra training compute, and the learned efficiency cannot transfer across architectures.
%
These costs motivate training-free designs on frozen models.
%
\textbf{Training-free Designs.}
Training-free methods cut cost on a frozen model via \textit{input-side} pruning or \textit{cache-side} compression.
%
On the \textit{input side}, AQuaUI~\citep{li2026aquaui} partitions screenshots with adaptive quadtrees, and related methods prune high-resolution screens or historical frames by spatio-temporal cues~\citep{xu2026guipruner,histprune2026}.
%
On the \textit{cache side}, GUI-KV~\citep{huang2025guikv} combines spatial saliency with temporal redundancy scoring.
%
ST-Lite~\citep{zhou2026stlite} couples component-centric saliency with trajectory-aware gating.
%
Despite these improvements, keeps are re-scored at every step, so adjacent keeps need not nest and a retired frame must be re-encoded or re-ranked.
%
In contrast, \method{} admits a nested keep once and contracts it as monotone session state efficiently.
\subsection{Visual Token Pruning}
Visual token pruning accelerates MLLM inference by removing redundant visual tokens, and existing methods either discard native tokens directly or aggregate them into synthetic ones.
%
\textbf{Pruning-based Methods.}
FastV~\citep{chen2024fastv} ranks visual tokens by attention statistics inside the language model.
%
DivPrune~\citep{alvar2025divprune} emphasizes feature dispersion to reduce redundancy.
%
CDPruner~\citep{zhang2025cdpruner} selects tokens that are both diverse and relevant to the instruction.
%
Nevertheless, selection discards tokens irreversibly, which motivates token merging as an alternative paradigm.
%
\textbf{Merging-based Methods.}
This paradigm aggregates discarded patches into compact synthetic summaries.
%
VisionZip~\citep{yang2025visionzip} merges visual tokens to extend the effective context length.
%
PruMerge+~\citep{shang2025prumerge} combines pruning with merging to adapt the visual token count.
%
Although effective, the selections of both paradigms drift across steps and demand rows that a contracted cache no longer holds.
%
Moreover, neither is tailored to GUI streams, where score-driven keeps collapse onto a few salient regions and sacrifice the spatial coverage that dense interfaces require.
%
In contrast, \method{} injects an interaction prior before prefill, repairs the spatial collapse at admission, and contracts the admitted keep into monotone session state for efficient serving.
